\documentclass[12pt]{report}

\usepackage{cmap}
\usepackage[T1,T2A]{fontenc}
\usepackage[utf8]{inputenc}
\usepackage[english, russian]{babel}
\usepackage{amssymb}
\usepackage{amsmath}
\usepackage{amsthm}
\usepackage{dsfont}
\usepackage{bm}
\usepackage{diagbox}
\usepackage[left=20mm,right=10mm,top=20mm,bottom=20mm,bindingoffset=2mm]{geometry}
\usepackage{indentfirst}
\usepackage[utf8]{inputenc}
\usepackage{float}
\usepackage[hidelinks]{hyperref}
\usepackage{graphicx}
\usepackage{xcolor}
\usepackage{listings}
\usepackage{minted}
\usepackage{booktabs}
\usepackage{array}
\usepackage{longtable}

\DeclareMathOperator{\N}{\mathbb{N}}
\DeclareMathOperator{\R}{\mathbb{R}}
\DeclareMathOperator{\Z}{\mathbb{Z}}
\DeclareMathOperator{\CC}{\mathbb{C}}
\DeclareMathOperator{\PP}{\mathrm{P}}
\DeclareMathOperator{\Expec}{\mathrm{E}}
\DeclareMathOperator{\Var}{\mathrm{Var}}
\DeclareMathOperator{\Cov}{\mathrm{Cov}}
\DeclareMathOperator{\asConv}{\xrightarrow{a.s.}}
\DeclareMathOperator{\LpConv}{\xrightarrow{Lp}}
\DeclareMathOperator{\pConv}{\xrightarrow{p}}
\DeclareMathOperator{\dConv}{\xrightarrow{d}}

\hypersetup{
	colorlinks=true,
	linkcolor=blue,
	citecolor=blue,
	urlcolor=blue
}

\lstset{language=Python, extendedchars=\true}

\lstdefinestyle{pythonstyle}{
	language=Python,
	backgroundcolor=\color{lightgray},
	commentstyle=\color{green},
	keywordstyle=\color{blue},
	stringstyle=\color{red},
	basicstyle=\ttfamily,
	frame=single,
	breaklines=true
}

\addto\captionsrussian{\renewcommand{\refname}{Список использованных источников}}

\begin{document}
	
	\begin{titlepage}
		\begin{center}
			\large{Федеральное государственное автономное образовательное учреждение высшего образования <<Национальный исследовательский университет ИТМО>>}
		\end{center}
		
		\vspace{15em}
		
		\begin{center}
			\huge{\textbf{Лабораторная работа №6}} \\
			\large{По дисциплине <<Вычислительная математика>>} \\
			\large{Вариант №12}
		\end{center}
		
		\vspace{5em}
		
		\begin{flushright}
			\textit{\large{Выполнил:}} \\
			\large{Студент группы P3218} \\
			\large{Михайлов Дмитрий} \\
			\large{Андреевич} \\
			\textit{\large{Преподаватель:}} \\
			\large{Малышева Татьяна} \\
			\large{Алексеевна}
		\end{flushright}
		
		\vspace{2cm}
		
		\begin{figure}[h]
			\centering
			\includegraphics[width=0.5\linewidth]{image.png}
		\end{figure}
		
		\begin{center}
			Санкт-Петербург \\
			2026 год
		\end{center}
	\end{titlepage}
	
	\tableofcontents
	\newpage
	
	\addcontentsline{toc}{section}{Цель работы}
	\section*{Цель работы}
	
	Целью лабораторной работы является изучение и практическая реализация численных методов решения задачи Коши для обыкновенных дифференциальных уравнений первого порядка.
	
	В рамках работы необходимо реализовать программу, позволяющую численно решать ОДУ при заданных начальных условиях, интервале интегрирования, шаге сетки и требуемой точности. Для варианта №12 согласно заданию используются следующие методы:
	\begin{itemize}
		\item метод Эйлера;
		\item метод Рунге--Кутта 4-го порядка;
		\item метод Милна.
	\end{itemize}
	
	Также целью работы является сравнение точности и особенностей указанных методов, анализ полученных результатов, оценка погрешности одношаговых методов по правилу Рунге и оценка точности многошагового метода по сравнению с точным решением задачи.
	
	\addcontentsline{toc}{section}{Блок-схемы используемых методов}
	\section*{Блок-схемы используемых методов}	
	
	\begin{figure}[H]
		\centering
		\includegraphics[width=0.3\linewidth]{Рисунок1.png}
		\caption{Блок-схема метода Эйлера}
	\end{figure}
	
	
	\addcontentsline{toc}{section}{Листинг программы}
	\section*{Листинг программы}
	\href{https://github.com/mysticslippers/math_comp_archive/blob/main/Lab6_comp_math/code/main.py}{Ссылка на репозиторий с кодом.}
	
	\addcontentsline{toc}{section}{Результат выполнения программы}
	\section*{Результат выполнения программы}
	
	
	\begin{figure}[H]
		\centering
		\includegraphics[width=0.6\linewidth]{screen1.png}
		\caption{Пример выполнения программы.}
	\end{figure}
	
	\begin{figure}[H]
		\centering
		\includegraphics[width=0.6\linewidth]{screen2.png}
		\caption{Пример выполнения программы.}
	\end{figure}
	
	\begin{figure}[H]
		\centering
		\includegraphics[width=0.8\linewidth]{graph.png}
		\caption{Вывод графиков функций.}
	\end{figure}
	
	
	\addcontentsline{toc}{section}{Вывод}
	\section*{Вывод}
	
	\noindent
	В ходе выполнения лабораторной работы была разработана программа для численного решения задачи Коши для обыкновенных дифференциальных уравнений первого порядка.
	
	В соответствии с вариантом №12 были реализованы три метода: метод Эйлера, метод Рунге--Кутта 4-го порядка и метод Милна. Для каждого метода были получены таблицы значений решения на заданной сетке, выполнена оценка точности и построены графики точного и приближенного решений.
	
	По результатам вычислительных экспериментов можно сделать следующие выводы. Метод Эйлера является самым простым по алгоритму и реализации, однако обладает наименьшей точностью среди рассмотренных методов. Его погрешность заметно возрастает при увеличении шага, поэтому для получения приемлемого результата требуется достаточно мелкая сетка.
	
	Метод Рунге--Кутта 4-го порядка показал значительно более высокую точность. Несмотря на то, что на каждом шаге необходимо вычислять несколько промежуточных коэффициентов, этот метод обеспечивает качественное приближение даже при сравнительно более крупном шаге. В рамках лабораторной работы именно данный метод показал наиболее устойчивые и точные результаты среди одношаговых методов.
	
	Метод Милна, являющийся многошаговым методом предиктор--корректор, также продемонстрировал высокую точность. Его особенностью является необходимость получения нескольких начальных точек другим методом, однако после этого вычисления выполняются эффективно. Результаты метода Милна оказались близки к точному решению, что подтверждает его практическую пригодность для численного решения задачи Коши.
	
	Таким образом, поставленная цель работы была достигнута: были изучены, реализованы и исследованы различные численные методы решения ОДУ. Работа показала, что выбор метода существенно влияет на точность результата, а использование методов более высокого порядка позволяет значительно уменьшить вычислительную погрешность.
	
	
\end{document}
